\section{第四章总习题}
\begin{problem}
    设函数$f$在$(a,b)$上连续,且$f(a+0)$与$f(b-0)$为有限值.证明:
    \begin{enumerate}
        \item $f$在$(a,b)$上有界;
        \item 若存在$\xi\in(a,b)$,使得$f(\xi)\geq\max\{f(a+0),f(b-0)\}$,则$f$在$(a,b)$上能取到最大值;
        \item $f$在$(a,b)$上一致连续.
    \end{enumerate}
\end{problem}
\begin{proof}
    \begin{enumerate}
        \item 令$f(a)=f(a+0),f(b)=f(b-0)$,于是$f(x)$在$[a,b]$上连续,由有界性定理,$f(x)$在$[a,b]$上有界,从而$f(x)$在$(a,b)$上有界;
        \item 由最大最小值定理,$f(x)$在$(a,b)$上存在最大值$M$.从而$M\geq f(\xi)\geq\max\{f(a+0),f(b-0)\}=\max\{f(a),f(b)\}$.若$M=f(\xi)$,而$\xi\in(a,b)$,因此$M$一定能在$(a,b)$内取到;若$M>f(\xi)$,则$M>f(a)$且$M>f(b)$,说明$M$不可能是端点值,从而$f(x)$能在$(a,b)$内取到最大值.
        \item $f$在$[a,b]$上连续,由康托定理,$f$在$[a,b]$上一致连续,从而$f$在$(a,b)$上一致连续.
    \end{enumerate}
\end{proof}
\begin{problem}
    设函数$f$在$(a,b)$上连续,且$f(a+0)=f(b-0)=+\infty$.证明:$f$在$(a,b)$上能取到最小值.
\end{problem}
\begin{proof}
    任取$x_0\in(a,b)$,$f(a+0)=+\infty$,由保号性,$\displaystyle\exists\delta_1\in(0,\dfrac{b-a}{2}),\forall x\in(a,a+\delta_1),f(x)\geq f(x_0)\Rightarrow x\in[a+\delta_1,b)$;$f(b-0)=+\infty$,由保号性,$\displaystyle\exists\delta_2\in(0,\dfrac{b-a}{2}),\forall x\in(b-\delta_2,b),f(x)\geq f(x_0)\Rightarrow x\in(a,b-\delta_2]$.当$x\in[a+\delta_1,b-\delta_2]$时,由最大最小值定理,$f(x)$存在最小值$f(c)$.当$x\in(a,a+\delta_1)\cup(b-\delta_2,b)$时,有$f(c)\leq f(x_0)\leq f(x)$.综上,$f$在$(a,b)$上能取到最小值.
\end{proof}